\section{Two Envelopes, Also known as the exchange paradox (Probability)}
There are two envelopes, one contains twice as much money as the other. You pick one at random and find $x$ dollars inside. What is the "expected value" of the amount in the other envelope? Given the option, should you switch? Does it matter how much you found in that first envelope? Does it matter if you open the envelope?

\subsection{Solution}
This is a strange paradox because it appears that the expected value is the average of $2x$ and $1/2x$ which are equally probable. The result is therefore $1.25x$ so you should switch. However, this would be true for any value of $x$ you find, so it wouldn't even matter if you open the envelope or not. This same logic suggests that you should also immediately switch back to the original envelope again, which doesn't make sense at all.

The fault is in the wrong working out of the expectation. The $x$ in each option is not the same $x$ since the sums $1/2x$, $x$ and $2x$ do not occur simultaneously. Let's do the calculation more carefully. We will denote the envelopes with the sums $1/2x$ and $x$. With probability of $1/2$ you got the envelope with $1/2x$ which means you $1/2x$ to gain if you switch. With the same probability of $1/2$ you got the envelope with $x$ in which case you have $1/2x$ to lose if you switch. This means that the expected value of switching is $0$ as it should be. No paradox at all. If we had denoted the two envelopes $x$ and $2x$ we would have got the same result: with probability $1/2$ we got the envelope with $x$ in it and can gain $x$ if we switch. With the same probability we got the envelope with $2x$ and stand to lose $x$ if we switch. The overall expected value is again $0$ as it should be.
